\chapter{Testing and Validation}
\label{ch:testing}

\section{K-Prototypes Results}
\label{ch:results_clustering}

The objective of clustering is to partition the patients into distinct groups based on their clinical profiles so that each group can be used as a scope for pattern mining. After removing outliers using the IQR criterion described in Section~\ref{ch:z_score}, exactly 21,995 (out of a total of 22,284) patients remained for the k-prototypes clustering. Patients marked as outliers (\texttt{is\_outlier = TRUE}) are included in the analysis but not included in the model fitting phase.

\subsubsection*{Benchmark Protocol}

In order to decide a final configuration, we benchmarked \textbf{63 feature-engineering versions} (\texttt{v1}\,$\ldots$\,\texttt{v63}), each picking a different subset of the 11 numerical and 10 categorical features built in Section~\ref{ch:feature_construction}. For every version, we performed a parameter sweep across the five k-prototypes hyperparameters: \texttt{scaler}~$\in$~\{\texttt{minmax}, \texttt{robust}\}, \texttt{gamma}~$\in$~\{\texttt{auto}, 0.07, 0.13, 0.17\}, \texttt{random\_state}~$\in$~\{7, 42, 100\}, \texttt{threshold}~$\in$~\{4, 6\}, and \texttt{k}~$\in$~\{4, 5, 6, 7\}. Every run generated a file named \texttt{conclusions.txt}, which included standard metrics (Silhouette, Gower Silhouette, Calinski-Harabasz, projected FAMD-Calinski-Harabasz, Davies-Bouldin, projected FAMD-Davies-Bouldin), cluster-wise Gower silhouette, and feature importance values (numerical $\eta^2$, categorical Cramer's V). Each iteration was ranked within its version, and the optimal iteration for each version was aggregated into a single \texttt{general\_conclusion.xlsx} sheet for ranking.

\subsubsection*{General Score and the Primacy of Gower}

The composite score that determines the ranking is the normalized sum of the aforementioned six individual scores, but it gives more weight to the \textbf{Gower Silhouette}, as explained in Section~\ref{ch:gower}. As for the type of input space, we have the mixed scenario (5~numerical and 8 categorical variables), while Silhouette, Calinski-Harabasz, and Davies-Bouldin indices assume the Euclidean metric; applying them to the one-hot encoding of categoricals results in an unwarranted emphasis on the numerical spread and masks the categorical nature of the data. However, the Gower distance first scales all numerical variables by their ranges and then assesses the dissimilarity between categorical values with binary agreement, just like k-prototypes does. The 2-D projection variants of Calinski-Harabasz and Davies-Bouldin are evaluated exclusively as tie-breakers because they are based on the dimensionality-reduced representation.

\subsubsection*{Feature Tuning per $k$}

After calculating the final scores, the upper end of the ranking was examined for each $k$, in order to select, among every possible cardinality, the feature set that obtained the most robust partitioning. Three optimal configurations were selected, which are listed in Table~\ref{tab:ch6_per_k_winners}. Scaler, threshold, and the FAMD measures are not listed, since the three winning sets share the same \texttt{minmax} scaler, threshold equal to \texttt{6}, and FAMD-DB ranging from 0.72 to 0.83.

\begin{table}[H]
\centering
\footnotesize
\renewcommand{\arraystretch}{1.1}
\begin{tabular}{@{}l c c c c c c c@{}}
\toprule
Config & Features ($n+c$) & $k$ & $\gamma$ & seed & Silh. & \textbf{Gower} & DB \\
\midrule
\textbf{v55} (final pick) & 5 + 8 & \textbf{4} & auto ($\approx$0.084) & 42 & 0.132 & \textbf{0.350} & 2.51 \\
v40 (best $k{=}5$)        & 4 + 8 & 5            & 0.17                  & 42 & 0.095 & 0.356          & 3.11 \\
v60 (best $k{=}6$)        & 4 + 8 & 6            & 0.17                  & 7  & 0.029 & 0.313          & 2.59 \\
\midrule
\multicolumn{8}{l}{\textit{Per-cluster Gower silhouette}} \\
\textbf{v55} ($k{=}4$) & \multicolumn{7}{l}{$c_0{=}0.423$,\quad $c_1{=}0.199$,\quad $c_2{=}0.320$,\quad $c_3{=}0.396$} \\
v40 ($k{=}5$)         & \multicolumn{7}{l}{$c_0{=}0.332$,\quad $c_1{=}0.537$,\quad $c_2{=}0.242$,\quad $c_3{=}0.203$,\quad $c_4{=}0.314$} \\
v60 ($k{=}6$)         & \multicolumn{7}{l}{$c_0{=}0.132$,\quad $c_1{=}0.242$,\quad $c_2{=}0.172$,\quad $c_3{=}0.331$,\quad $c_4{=}0.483$,\quad $c_5{=}0.402$} \\
\midrule
% Titlurile coloanelor pentru features
\multicolumn{3}{@{}l}{\textbf{v55}} & \multicolumn{3}{l}{v40} & \multicolumn{2}{l@{}}{v60} \\
\cmidrule(r){1-3} \cmidrule(lr){4-6} \cmidrule(l){7-8}

\multicolumn{8}{l}{\textit{Numerical Features}} \\
\multicolumn{3}{@{}l}{\begin{tabular}[t]{@{}l@{}}
\texttt{age} \\ \texttt{active\_conditions\_count} \\ \texttt{total\_medications} \\ \texttt{total\_procedures} \\ \texttt{preventive\_to\_reactive\_ratio}
\end{tabular}} &
\multicolumn{3}{l}{\begin{tabular}[t]{@{}l@{}}
\texttt{age} \\ \texttt{active\_conditions\_count} \\ \texttt{total\_medications} \\ \texttt{lifetime\_medication\_days}
\end{tabular}} &
\multicolumn{2}{l@{}}{\begin{tabular}[t]{@{}l@{}}
\texttt{age} \\ \texttt{active\_conditions\_count} \\ \texttt{total\_medications} \\ \texttt{preventive\_to\_reactive\_ratio}
\end{tabular}} \\

\addlinespace % Adaugă puțin spațiu între blocul numeric și cel categoric

\multicolumn{8}{l}{\textit{Categorical Features}} \\
\multicolumn{3}{@{}l}{\begin{tabular}[t]{@{}l@{}}
\texttt{gender} \\ \texttt{pregnancy\_history} \\ \texttt{has\_heart\_disease} \\ \texttt{has\_diabetes} \\ \texttt{has\_hypertension} \\ \texttt{condition\_burden} \\ \texttt{bmi\_type} \\ \texttt{dominant\_careplan}
\end{tabular}} &
\multicolumn{3}{l}{\begin{tabular}[t]{@{}l@{}}
\texttt{gender} \\ \texttt{pregnancy\_history} \\ \texttt{has\_heart\_disease} \\ \texttt{has\_diabetes} \\ \texttt{has\_hypertension} \\ \texttt{condition\_burden} \\ \texttt{bmi\_type} \\ \texttt{cholesterol\_risk\_category}
\end{tabular}} &
\multicolumn{2}{l@{}}{\begin{tabular}[t]{@{}l@{}}
\texttt{gender} \\ \texttt{pregnancy\_history} \\ \texttt{has\_heart\_disease} \\ \texttt{has\_diabetes} \\ \texttt{has\_hypertension} \\ \texttt{condition\_burden} \\ \texttt{bmi\_type} \\ \texttt{cholesterol\_risk\_category}
\end{tabular}} \\
\bottomrule
\end{tabular}
\caption{Winning configurations per $k$. Base feature sets are nearly identical across the three winners, with minor substitutions in secondary numerical metrics and categorical variables (e.g., \texttt{v55} utilizes \texttt{dominant\_careplan} instead of \texttt{cholesterol\_risk\_category}).}
\label{tab:ch6_per_k_winners}
\end{table}

\subsubsection*{Final Configuration and Rationale}

\texttt{v55\_sminmax\_gauto\_rs42\_th6\_k4}, despite having a slightly lower overall Gower silhouette score of 0.350 than \texttt{v40-$k{=}5$} (0.356), has \textbf{four} clusters that all essentially meet or exceed the heuristic value of 0.20, indicating stability (\textit{min $=$~0.199, max $=$~0.423}). For the optimum with five clusters, there is only \textbf{one} cluster that has a score of 0.203, barely exceeding the heuristic threshold. As for the optimum with six clusters, there are \textbf{three} clusters with scores below the heuristic value. It is essential to select such strict criteria since the clustering process will be used in the next step of PrefixSpan mining. The best $k$ value would therefore be the smallest value for which all clusters remain structurally cohesive: \texttt{active\_conditions\_count} ($\eta^2{=}0.74$) and \texttt{age} ($\eta^2{=}0.66$) dominate the numerical contribution, while \texttt{condition\_burden} (Cramer's V $=$~0.80), \texttt{gender} (0.70), \texttt{has\_diabetes} (0.69), and \texttt{has\_heart\_disease} (0.66) dominate the categorical side.

\subsubsection*{Visual Inspection of the Final Partition}

\begin{figure}[H]
\centering
\includegraphics[width=1\linewidth]{figs/chapter_6/clustering_famd_scatter_plot.png}
\caption{FAMD 2-D projection of the four clusters (\texttt{v55}, $k{=}4$).}
\label{fig:ch6_clu_famd}
\end{figure}

The scatterplot of FAMD reveals how the four clusters exhibit clear decision boundaries along the first two components of FAMD, although k-prototypes never uses such an embedding natively.

\begin{figure}[H]
\centering
\includegraphics[width=1\linewidth]{figs/chapter_6/clustering_heatmap.png}
\caption{Heatmap of z-scored numerical feature means per cluster.}
\label{fig:ch6_clu_heatmap}
\end{figure}

The heatmap indicates an obvious morbidity gradient with respect to cluster~0 representing younger patients with low morbidity and cluster~3 at the other extreme with high morbidity. Clusters~1 and ~2 represent intermediate states. However, the skew in \texttt{preventive\_to\_reactive\_ratio} is different for both clusters. In this visualization, warmer colors (red) indicate values that are higher than the average while blue colors reflect values lower than the dataset average.

\begin{figure}[H]
\centering
\includegraphics[width=1\linewidth]{figs/chapter_6/clustering_numerical_graphs.png}
\caption{Per-cluster means of the five numerical features.}
\label{fig:ch6_clu_numerical}
\end{figure}

Examining the per-cluster bar plot of the five numerical features, we can clearly see the importance of \texttt{active\_conditions\_count} and \texttt{age}. These two features primarily drive the divergence between clusters~2 and~3 from clusters~0 and~1.

\begin{figure}[H]
\centering
\includegraphics[width=1\linewidth]{figs/chapter_6/clustering_categorical_graphs.png}
\caption{Per-cluster proportions of the eight categorical features.}
\label{fig:ch6_clu_categorical}
\end{figure}

The categorical panel corroborates these findings from the other side: \texttt{condition\_burden} cleanly stratifies the four clusters, \texttt{has\_diabetes} and \texttt{has\_heart\_disease} concentrate inside the high-burden clusters, and \texttt{gender}~/~\texttt{pregnancy\_history} account for the primary variance separating cluster~1 (women of reproductive age) from the rest.

\label{ch:cluster_results}

\section{Sequence Mining Results}
\label{ch:results_sequences}

PrefixSpan (Chapter~\ref{ch:prefix_span_theory}) was executed per scope; a separate run for each of the four clusters in Section~\ref{ch:cluster_results} led to a distinct set of patterns, while the cluster of outliers (\texttt{cluster\_id = -1}) was processed separately due to its special status of containing only individuals with a high-acuity signature, having been excluded during k-prototypes clustering. The resulting number of frequent patterns varied between 175 and 252 depending on the scope. In what follows, we present the \textbf{top 10} for each scope based on geometric mean lift (\texttt{gml}). This metric favors patterns whose constituent items are specifically frequent within the target scope, rather than simply common across the entire dataset. Each table below includes support size (\texttt{pts}), relative support within the scope (\texttt{\%}), \texttt{gml} score, and an ordered itemset with the notation \texttt{C:} (condition), \texttt{M:} (medication), \texttt{P:} (procedure), \texttt{CP:} (careplan), and \texttt{[dental]} when the pattern was identified as dentistry-specific during post-filtering. Generic medication names are abbreviated to INN form.

\begin{table}[H]
\centering
\footnotesize
\renewcommand{\arraystretch}{1.2}
\begin{tabular}{@{}r r r r p{10.5cm}@{}}
\toprule
\textbf{\#} & \textbf{pts} & \textbf{\%} & \textbf{gml} & \textbf{Pattern} \\
\midrule
1  &  252 &  3.6 & 0.97 & C:Laceration $\to$ CP:Wound care $\to$ P:Suture open wound \\
2  &  113 &  1.6 & 0.78 & C:Fracture of bone $\to$ CP:Fracture care $\to$ P:Bone immobilization \\
3  &  106 &  1.5 & 0.96 & M:Fluticasone/salmeterol DPI $\to$ M:albuterol $\to$ P:Spirometry \\
4  &  432 &  6.1 & 0.70 & C:Acute bronchitis $\to$ CP:Respiratory therapy $\to$ M:Acetaminophen \\
5  & 2766 & 39.0 & 0.82 & [dental] M:sodium fluoride gel $\to$ P:Dental fluoride treatment $\to$ P:Oral health education \\
6  &  295 &  4.2 & 0.62 & C:Acute bronchitis $\to$ CP:Respiratory therapy $\to$ P:Measurement of respiratory function \\
7  & 2705 & 38.1 & 0.73 & [dental] P:Examination of gingivae $\to$ M:sodium fluoride $\to$ P:Dental fluoride $\to$ P:Oral health education \\
8  &  127 &  1.8 & 0.73 & C:Sprain $\to$ C:Sprain of ankle $\to$ CP:Physiotherapy care plan \\
9  &  229 &  3.2 & 0.44 & C:Acute bronchitis $\to$ CP:Respiratory therapy $\to$ P:Measurement resp. function $\to$ M:Acetaminophen \\
10 &  104 &  1.5 & 0.47 & C:Streptococcal sore throat $\to$ M:Penicillin V $\to$ P:Throat culture \\
\bottomrule
\end{tabular}
\caption{Cluster~0 (young / low-burden) - top 10 patterns out of 238.}
\label{tab:ch6_seq_c0}
\end{table}

\begin{table}[H]
\centering
\footnotesize
\renewcommand{\arraystretch}{1.2}
\begin{tabular}{@{}r r r r p{10.5cm}@{}}
\toprule
\textbf{\#} & \textbf{pts} & \textbf{\%} & \textbf{gml} & \textbf{Pattern} \\
\midrule
1  &  154 &  2.4 & 0.97 & C:Laceration $\to$ CP:Wound care $\to$ P:Suture open wound \\
2  &  121 &  1.9 & 0.92 & C:Sterilization requested $\to$ P:Sterilization education $\to$ C:Hist. tubal ligation $\to$ CP:Minor surgery $\to$ P:Postoperative care \\
3  &  123 &  1.9 & 0.90 & C:History of tubal ligation $\to$ CP:Minor surgery care $\to$ P:Postoperative care \\
4  &  107 &  1.7 & 0.81 & P:Sterilization ed. $\to$ C:Hist. tubal ligation $\to$ CP:Minor surgery $\to$ P:Ligation of fallopian tubes $\to$ P:Postoperative care \\
5  &  147 &  2.3 & 0.96 & M:Fluticasone/salmeterol DPI $\to$ M:albuterol $\to$ P:Spirometry \\
6  &   99 &  1.5 & 0.96 & P:Prone position $\to$ M:Furosemide $\to$ P:Oxygen by mask $\to$ P:Prone position $\to$ M:Furosemide \\
7  &  308 &  4.8 & 0.71 & C:Acute bronchitis $\to$ CP:Respiratory therapy $\to$ M:Acetaminophen \\
8  & 2058 & 32.1 & 0.82 & [dental] M:sodium fluoride $\to$ P:Dental fluoride $\to$ P:Oral health education \\
9  &  109 &  1.7 & 0.79 & CP:Minor surgery $\to$ P:Ligation of fallopian tubes $\to$ P:Postoperative care \\
10 &  432 &  6.7 & 0.72 & C:Normal pregnancy $\to$ P:Standard pregnancy test $\to$ P:Ultrasound fetal viability \\
\bottomrule
\end{tabular}
\caption{Cluster~1 (women of reproductive age) - top 10 patterns out of 252.}
\label{tab:ch6_seq_c1}
\end{table}

\begin{table}[H]
\centering
\footnotesize
\renewcommand{\arraystretch}{1.2}
\begin{tabular}{@{}r r r r p{10.5cm}@{}}
\toprule
\textbf{\#} & \textbf{pts} & \textbf{\%} & \textbf{gml} & \textbf{Pattern} \\
\midrule
1  &  154 &  4.5 & 0.95 & M:Fluticasone/salmeterol DPI $\to$ M:albuterol $\to$ P:Spirometry \\
2  &  164 &  4.8 & 0.89 & M:Furosemide $\to$ P:Oxygen by mask $\to$ P:Prone position $\to$ M:Furosemide \\
3  &  138 &  4.0 & 0.81 & M:Epoetin Alfa $\to$ P:Renal dialysis $\to$ M:Epoetin Alfa \\
4  & 1222 & 35.6 & 0.80 & [dental] M:sodium fluoride $\to$ P:Dental fluoride $\to$ P:Oral health education \\
5  & 1163 & 33.9 & 0.77 & [dental] P:Examination of gingivae $\to$ M:sodium fluoride $\to$ P:Dental fluoride \\
6  &  105 &  3.1 & 0.74 & M:albuterol $\to$ P:Pulmonary rehabilitation $\to$ P:Spirometry \\
7  &  101 &  2.9 & 0.70 & M:Fluticasone/salmeterol DPI $\to$ M:albuterol $\to$ P:Pulmonary rehab $\to$ P:Spirometry \\
8  & 1162 & 33.8 & 0.69 & [dental] P:Examination gingivae $\to$ M:sodium fluoride $\to$ P:Dental fluoride $\to$ P:Oral health education \\
9  & 1111 & 32.3 & 0.68 & [dental] P:Removal subgingival plaque $\to$ P:Examination gingivae $\to$ M:sodium fluoride $\to$ P:Dental fluoride \\
10 &   97 &  2.8 & 0.67 & P:Prone position $\to$ M:Furosemide $\to$ P:Heart failure education \\
\bottomrule
\end{tabular}
\caption{Cluster~2 (hypertensive adults, polypharmacy) - top 10 patterns out of 202.}
\label{tab:ch6_seq_c2}
\end{table}

\begin{table}[H]
\centering
\footnotesize
\renewcommand{\arraystretch}{1.2}
\begin{tabular}{@{}r r r r p{10.5cm}@{}}
\toprule
\textbf{\#} & \textbf{pts} & \textbf{\%} & \textbf{gml} & \textbf{Pattern} \\
\midrule
1  &   73 &  1.6 & 1.00 & C:Acute allergic reaction $\to$ P:Injection of epinephrine $\to$ M:prednisone \\
2  &  137 &  2.9 & 0.98 & C:Laceration $\to$ CP:Wound care $\to$ P:Suture open wound \\
3  &  134 &  2.9 & 0.84 & CP:Self-care interventions $\to$ P:Allergy screening test $\to$ M:Epinephrine auto-injector \\
4  &  169 &  3.6 & 0.81 & C:Fracture of bone $\to$ CP:Fracture care $\to$ P:Bone immobilization \\
5  &  459 &  9.8 & 0.78 & C:Acute bronchitis $\to$ CP:Respiratory therapy $\to$ M:Acetaminophen \\
6  &   48 &  1.0 & 0.91 & C:Fracture of clavicle $\to$ P:Admission to orthopedic dept. $\to$ P:Plain X-ray of clavicle \\
7  &   93 &  2.0 & 0.91 & C:History of seizure $\to$ C:Seizure disorder $\to$ P:Electroencephalogram \\
8  & 1873 & 40.1 & 0.83 & [dental] M:sodium fluoride $\to$ P:Dental fluoride $\to$ P:Oral health education \\
9  & 1856 & 39.8 & 0.83 & [dental] P:Examination of gingivae $\to$ M:sodium fluoride $\to$ P:Dental fluoride \\
10 &   87 &  1.9 & 0.80 & C:Concussion injury $\to$ C:Concussion no LOC $\to$ CP:Head injury rehabilitation \\
\bottomrule
\end{tabular}
\caption{Cluster~3 (pediatric) - top 10 patterns out of 175.}
\label{tab:ch6_seq_c3}
\end{table}

\begin{table}[H]
\centering
\footnotesize
\renewcommand{\arraystretch}{1.2}
\begin{tabular}{@{}r r r r p{10.5cm}@{}}
\toprule
\textbf{\#} & \textbf{pts} & \textbf{\%} & \textbf{gml} & \textbf{Pattern} \\
\midrule
1  &    7 &  2.5 & 1.00 & P:Oxygen by mask $\to$ P:Prone position $\to$ M:Furosemide $\to$ P:Oxygen by mask $\to$ P:Prone $\to$ M:Furosemide \\
2  &    6 &  2.1 & 1.00 & M:Fluticasone/salmeterol DPI $\to$ M:albuterol $\to$ P:Pulmonary rehab $\to$ P:Spirometry \\
3  &  142 & 49.8 & 0.97 & P:Combined chemotherapy + radiation $\to$ M:Cisplatin $\to$ P:Combined chemotherapy + radiation \\
4  &    9 &  3.2 & 0.90 & P:Renal disorder medication review $\to$ M:tacrolimus [Prograf] $\to$ P:Transplantation of kidney \\
5  &  116 & 40.7 & 0.85 & M:Cisplatin $\to$ M:Paclitaxel $\to$ P:Combined chemotherapy + radiation \\
6  &  132 & 46.3 & 0.85 & M:Cisplatin $\to$ P:Combined chemotherapy + radiation $\to$ P:Hospice care \\
7  &   31 & 10.9 & 0.78 & P:Renal dialysis $\to$ M:Epoetin Alfa $\to$ P:Renal dialysis $\to$ M:Epoetin Alfa \\
8  &   35 & 12.3 & 0.76 & [dental] M:sodium fluoride $\to$ P:Dental fluoride $\to$ P:Oral health education \\
9  &  105 & 36.8 & 0.72 & P:Combined chemo+radiation $\to$ M:Cisplatin $\to$ M:Paclitaxel $\to$ P:Combined chemo+radiation $\to$ P:Hospice \\
10 &   81 & 28.4 & 0.67 & P:Combined chemo+radiation $\to$ P:Hearing examination $\to$ M:Cisplatin $\to$ P:Combined chemo+radiation \\
\bottomrule
\end{tabular}
\caption{Cluster -1 (oncology / renal / heart failure) - top 10 patterns out of 224.}
\label{tab:ch6_seq_outliers}
\end{table}

\subsection{Per-Cluster Reading}

Looking across the five tables (\textbf{Tables~\ref{tab:ch6_seq_c0} through~\ref{tab:ch6_seq_outliers}}), each scope produces a clinically distinct dominant signal:

\begin{itemize}
    \item \textbf{Cluster~0 (young / low-burden).} As illustrated by Table~\ref{tab:ch6_seq_c0}, the most prominent trends consist of episodic acute care: laceration $\rightarrow$ suture (rank~1), fracture $\rightarrow$ immobilization (rank~2), acute bronchitis $\rightarrow$ acetaminophen (rank~4), and the ubiquitous dental fluoride / gingivitis combination. No chronic-disease chain rises to the top, consistent with the cluster's near-zero \texttt{active\_conditions\_count} from Section~\ref{ch:results_clustering}.

    \item \textbf{Cluster~1 (reproductive women).} Refer to Table~\ref{tab:ch6_seq_c1}. The discriminating features of this cluster include reproductive health care pathways, such as tubal ligation sequences (ranks 2-4, gml 0.81-0.92) and the sequence of normal pregnancy followed by a fetal viability ultrasound (rank 10).

    \item \textbf{Cluster~2 (adults with hypertension, use of multiple medications).} This cluster is clearly seen in Table~\ref{tab:ch6_seq_c2} as one characterized by chronic disease management, addressing fluid overload via furosemide / prone positioning / oxygen sequences (rank~2, rank~10), the renal dialysis / Epoetin Alfa loop (rank~3), and obstructive airway pulmonary rehabilitation (ranks~6-7). Further down the extended list (ranks~14-15) appears the sequence of lisinopril with depression and substance use screening.

    \item \textbf{Cluster~3 (pediatrics).} Table~\ref{tab:ch6_seq_c3} demonstrates a prevalence of trauma and allergic events: acute allergic reaction $\rightarrow$ epinephrine $\rightarrow$ prednisone (rank~1, gml 1.00), epinephrine auto-injector instruction (rank~3), fracture chains for clavicle, arm, ankle, bone (ranks~4, 6), seizures/EEG (rank~7), and concussion rehabilitation (rank~10). Otitis-media antibiotic chains appear further down the long tail (rank 27+ in the source file).

    \item \textbf{Outliers (cluster~$-1$).} These are the high-acuity pathways that the IQR filter has isolated from the clustering process and now manifest in an isolated form (Table~\ref{tab:ch6_seq_outliers}), with combined chemotherapy and radiation treatment involving cisplatin/paclitaxel topping the list (ranks~3, 5, 6, 9-10, and more from the elongated tail), kidney dialysis combined with Epoetin Alfa (rank~7), kidney transplantation managed with tacrolimus (rank~4), and a heart-failure pathway of furosemide/prone position/oxygen combination (rank~1). Considering the outliers as an independent set is the only reason why these expensive but rare pathways have appeared at all. They fall below the minimum-support threshold when mined against the global population.
\end{itemize}

\vspace{1em}
As a whole, the five pattern sets demonstrate the semantic separability of the four-cluster partition: each cluster’s top patterns characterize a distinct patient profile, which is required to enable cluster-specific pathway analytics in the dashboard.